\documentclass[10pt,a4paper]{amsart}
\usepackage[margin=19mm]{geometry}
\usepackage{amsmath,amssymb,mathtools}
\usepackage[scr=rsfs,cal=euler]{mathalfa}
\usepackage{xfrac}
\usepackage[unicode,hidelinks,bookmarks=false]{hyperref}
\newcommand{\ie}{i.e.\ }
\newcommand{\cf}{cf.\ }
\newcommand{\resp}{respectively\ }
\newcommand{\Subsection}{\subsection}
\providecommand{\nobreakdash}{\nobreak\hbox{-}}
\setlength{\emergencystretch}{2em}
\multlinegap0pt
\usepackage{bm}
\usepackage{bbm}
\usepackage{dsfont}
\usepackage{xspace}
\usepackage{cancel}
\usepackage{mathtools}
\usepackage{amsthm}
\makeatletter
\@ifundefined{thm}{\newtheorem{thm}{Thm}}{}
\@ifundefined{lem}{\newtheorem{lem}[thm]{Lemma}}{}
\@ifundefined{thm*}{\newtheorem*{thm*}{Theorem}}{}
\makeatother
\newcommand{\diag}{\operatorname{diag}}
\newcommand{\per}{\operatorname{per}}
\newcommand{\sym}{\operatorname{sym}}
\title[Two Proofs of the Hamiltonian Cycle Identity]{Two Proofs of the Hamiltonian Cycle Identity}
\author{Hamilton Sawczuk, Edinah Gnang}
\date{Source: arXiv:2510.02473v1 (CC BY 4.0)}
\begin{document}
\maketitle

\noindent\emph{Source and licence.} Two Proofs of the Hamiltonian Cycle Identity. Version arXiv:2510.02473v1 (CC BY 4.0), distributed under the Creative Commons Attribution 4.0 International licence, as recorded in the arXiv licence metadata for this exact version and shown on the abstract page https://arxiv.org/abs/2510.02473v1. Full rights notice: licenses/arxiv-2510.02473-CC-BY-4.0.txt.\par\smallskip

\noindent\textit{Editorial scope.} Counted unit: the Thm* whose source label is \texttt{None} in the article (source0.tex lines 527--533, source label \texttt{None}), delivered together with the article's own complete original proof of it (source0.tex lines 535--616, one \texttt{proof} environment of 82 lines). No mathematics is written, repaired or re-derived: statement and proof are the article's own text line for line. Required context retained uncounted: lem:PHn (lines 451--456). Every definition, display and label of those lines is retained unchanged. Omitted from this package and not redistributed: 1--450, 457--526, 534--534, 617--691. Nothing inside a retained excerpt is omitted or altered: every retained body is delivered exactly as the article's lines are written, so no omission receipt of any kind accompanies this package. Citations: the retained excerpts carry no citation command at all, so no bibliography entry of the article is transcribed and no reference is redistributed. Reference adaptation: none was needed and none was made; every reference inside the delivered unit resolves inside it, so no unresolved reference or citation is produced. Printed slips kept literal: no printed word, symbol, hypothesis or number of the counted statement or of its proof was altered. Layout and class: the article's own class is \texttt{article} with packages including maa-monthly, tikz; the wrapper is the lane's pinned \texttt{amsart} editorial wrapper at 10pt on A4, which restates the theorem-like environments the retained text needs and adds the article's own macro definitions that the retained text needs (\texttt{\textbackslash diag}, \texttt{\textbackslash per}, \texttt{\textbackslash sym}), with extra wrapper packages (bm, bbm, dsfont, xspace, cancel, mathtools). The delivered numbering is the unit's own. Assets: no retained span contains a figure environment, a graphics-inclusion command, a PDF-page inclusion, a table or a TikZ picture; no image file is copied into this package, the package declares no asset dependency and no image file is redistributed with it. Licence: version arXiv:2510.02473v1 of this article is distributed under the Creative Commons Attribution 4.0 International licence, as recorded in the arXiv licence metadata for this exact version and shown on the abstract page https://arxiv.org/abs/2510.02473v1. A full rights notice is delivered at licenses/arxiv-2510.02473-CC-BY-4.0.txt. Characters: every retained excerpt is pure ASCII, so the delivered engine, XeLaTeX with its default Latin Modern text font, reports no missing character.
\begin{lem}\label{lem:PHn}
    The following expression is a symbolic listing of Hamiltonian cycles on $[n]$.
    \[
        P_{HC_n}(A)=\boldsymbol{\partial}_{[n]}\sum_{j\in[n]}a_{n,j}x_j\det(\diag(A\boldsymbol{x})-A\diag(\boldsymbol{x}))_{[n-1]}
    \]
\end{lem}

\begin{thm*}
    The Hamiltonian cycle polynomial satisfies the identity
    \[
        P_{HC_n}(A)=\sum_{S\subseteq [n-1]}\det(-A_S)\per(A_{[n]\setminus S}),
    \]
    where $A_S$ denotes the principle submatrix of $A$ indexed by $S$, and $\det(A_{\emptyset}):=1$.
\end{thm*}

\begin{proof}
    By Lemma \ref{lem:PHn}, we have
    \[
        P_{HC_n}(A)=\boldsymbol{\partial}_{[n]}\sum_{j\in[n]}a_{n,j}x_j\cdot\det(\diag(A\boldsymbol{x})-A\diag(\boldsymbol{x}))_{[n-1]}
    \]
    is a symbolic listing of Hamiltonian cycles. Now observe
    \[
        \det(\diag(A\boldsymbol{x})-A\diag(\boldsymbol{x}))_{[n-1]}=\det(-A\diag(\boldsymbol{x}) + \diag(A\boldsymbol{x}))_{[n-1]},
    \]
    and by the determinant sum lemma
    \[
        \det(-A\diag(\boldsymbol{x}) + \diag(A\boldsymbol{x}))_{[n-1]}
    \]
    \[
        =\sum_{S\subseteq[n-1]}\det(-A\diag(\boldsymbol{x}))_S\det(\diag(A\boldsymbol{x}))_{[n-1]\setminus S}
    \]
    \[
        =\sum_{S\subseteq[n-1]}\det(-A\diag(\boldsymbol{x}))_S\per(\diag(A\boldsymbol{x}))_{[n-1]\setminus S}
    \]
    since the determinant and permanent of a diagonal matrix are the same. Plugging this in, we find $P_{HC_n}$ equals
    \[
        \boldsymbol{\partial}_{[n]}\sum_{j\in[n]}a_{n,j}x_j\sum_{S\subseteq[n-1]}\det(-A\diag(\boldsymbol{x}))_S\per(\diag(A\boldsymbol{x}))_{[n-1]\setminus S},
    \]
    which can be factored as
    \[
        \sum_{S\subseteq[n-1]}\boldsymbol{\partial}_{[n]}\det(-A\diag(\boldsymbol{x}))_S\sum_{j\in[n]}a_{n,j}x_j\per(\diag(A\boldsymbol{x}))_{[n-1]\setminus S}.
    \]
    Now fix $S\subseteq[n-1]$ and consider a single summand. Observe that the left and right hand factors above have degree $|S|$ and $n-|S|$ respectively in the variables of $\boldsymbol{x}$. By the multivariable product rule, the summand equals
    \[
        \sum_{T\subseteq[n]}\boldsymbol{\partial}_{T}\det(-A\diag(\boldsymbol{x}))_S
        \cdot\boldsymbol{\partial}_{[n]\setminus T}\sum_{j\in[n]}a_{n,j}x_j\per(\diag(A\boldsymbol{x}))_{[n-1]\setminus S}
    \]
    The left hand factor vanishes unless $T\subseteq S$. But if $|T|<|S|$ then the right hand factor vanishes since it has total degree $n-|S|$ but a degree $n-|T|>n-|S|$ differential operator. Thus the only non-vanishing term corresponds to $T=S$, and $P_{HC_n}(A)$ equals
    \[
        \sum_{S\subset[n-1]}\boldsymbol{\partial}_{S}\det(-A\diag(\boldsymbol{x}))_S \cdot \boldsymbol{\partial}_{[n]\setminus S}\sum_{j\in[n]}a_{n,j}x_j\per(\diag(A\boldsymbol{x}))_{[n-1]\setminus S}.
    \]
    Now
    \[
        \boldsymbol{\partial}_{S}\det(A\diag(\boldsymbol{x}))_S=\det(A)_S\cdot\boldsymbol{\partial}_S\det(\diag(\boldsymbol{x}))_S
    \]
    \[
        =\det(A)_S\cdot\boldsymbol{\partial}_{S}\prod_{i\in S}x_i=\det(A)_S
    \]
    which implies
    \[
        P_{HC_n}=\sum_{S\subset[n-1]}\det(-A)_S\cdot\boldsymbol{\partial}_{[n]\setminus S}\sum_{j\in[n]}a_{n,j}x_j\per(\diag(A\boldsymbol{x}))_{[n-1]\setminus S}.
    \]
    Finally, we examine the right hand factor
    \[
        \boldsymbol{\partial}_{[n]\setminus S}\sum_{j\in[n]}a_{n,j}x_j\per(\diag(A\boldsymbol{x}))_{[n-1]\setminus S}
    \]
    \[
        =\boldsymbol{\partial}_{[n]\setminus S}\sum_{j\in[n]}a_{n,j}x_j\prod_{i\in[n-1]\setminus S}\left(\sum_{k\in[n]}a_{i,k}x_{k}\right).
    \]
    Since the total degree is $n-|S|$, any non-vanishing term must have degree one in the variables of $\boldsymbol{x}$ indexed by $[n]\setminus S$ and degree zero in those indexed by $S$. Thus we can restrict the above sums to $j,k\in[n]\setminus S$ to write
    \[
        =\boldsymbol{\partial}_{[n]\setminus S}\sum_{j\in[n]\setminus S}a_{n,j}x_j\prod_{i\in[n-1]\setminus S}\left(\sum_{k\in[n]\setminus S}a_{i,k}x_{k}\right)
    \]
    \[
        =\boldsymbol{\partial}_{[n]\setminus S}\left(\prod_{i\in[n-1]\setminus S}\sum_{k\in[n]\setminus S}a_{i,k}x_{k}\right)\left(\sum_{j\in[n]\setminus S}a_{n,j}x_j\right)
    \]
    \[
        =\boldsymbol{\partial}_{[n]\setminus S}\prod_{i\in[n]\setminus S}\sum_{k\in[n]\setminus S}a_{i,k}x_{k}.
    \]
    Now we switch the order of product and summation, yielding
    \[
        =\boldsymbol{\partial}_{[n]\setminus S}\sum_{f:([n]\setminus S)\rightarrow([n]\setminus S)}\prod_{i\in[n]\setminus S}a_{i,f(i)}x_{f(i)},
    \]
    but as before, only monomials with in-degree sequence $(1,1,\hdots,1)$, i.e. permutations, will remain after the action of partial derivative operator. Thus we restrict the sum to find
    \[
        =\boldsymbol{\partial}_{[n]\setminus S}\sum_{\sigma\in\sym([n]\setminus S)}\left(\prod_{i\in[n]\setminus S}a_{i,\sigma(i)}x_{\sigma(i)}\right) =\boldsymbol{\partial}_{[n]\setminus S}\per(A\diag(\boldsymbol{x}))_{[n]\setminus S}.
    \]
    As with the determinant, 
    \[
        \boldsymbol{\partial}_{S}\per(A\diag(\boldsymbol{x}))_S=\per(A)_S\cdot\boldsymbol{\partial}_{S}\prod_{i\in S}x_i=\per(A)_S,
    \]
    so the right hand factor is equal to $\per(A_{[n]\setminus S})$, and therefore
    \[
        P_{HC_n}(A)=\sum_{S\subseteq[n-1]}\det(-A_S)\per(A_{[n]\setminus S})
    \]
    as desired.
\end{proof}

\end{document}
